\chapter{INTRODUÇÃO}
A Introdução deve ser redigida em texto corrido, de forma objetiva, fluida e concisa. Não utilize tópicos, subtópicos ou listas com marcadores neste capítulo. Embora seja o primeiro capítulo do trabalho, recomenda-se que sua revisão final seja feita após a conclusão da Monografia.

A Introdução deve contextualizar o leitor sobre o tema e a pesquisa realizada. Sugere-se a seguinte estrutura lógica de parágrafos para a construção do seu texto:
\begin{uemaenum}
    \item \textbf{Contextualização e Tema:} Apresente o tema central da sua monografia. Defina os conceitos principais e situe o leitor sobre o cenário atual desse assunto (1º e 2º Parágrafos);
    \item \textbf{Problematização:} Apresente a situação-problema encontrada e culmine na sua pergunta norteadora (o problema de pesquisa em forma de interrogação) (3º Parágrafo);
    \item \textbf{Justificativa:} Explique a relevância do tema. Por que esta pesquisa é importante para a sociedade e para o ambiente acadêmico? (4º Parágrafo);
    \item \textbf{Objetivos:} Apresente o Objetivo Geral e os Objetivos Específicos do trabalho. Lembre-se: insira-os no meio do texto corrido utilizando conectivos (ex: ``Para tanto, este estudo tem como objetivo geral... e, especificamente, busca-se...''), sem listá-los em tópicos (5º Parágrafo);
    \item \textbf{Metodologia:} Resuma brevemente a Metodologia adotada (ex: se é uma pesquisa bibliográfica, estudo de caso, qualitativa, etc.) e as principais fontes utilizadas (6º Parágrafo);
    \item \textbf{Estrutura da Monografia:} Explique como o trabalho está organizado. Indique o que o leitor encontrará em cada um dos capítulos seguintes (ex: ``O presente trabalho está estruturado em quatro capítulos. O primeiro traz esta Introdução... O segundo capítulo aborda o Referencial Teórico sobre X...'') (7º e 8º Parágrafos).
\end{uemaenum}
